% source: https://github.com/pfradin/luadraw/discussions/288#discussioncomment-17604343 (pfradin, block 1)
\documentclass[border=3pt]{standalone}% compile with lualatex only
\usepackage[svgnames]{xcolor}
\usepackage[3d]{luadraw}%https://github.com/pfradin/luadraw
\usepackage{fourier-otf}
\begin{document}
\begin{luadraw}{name=cylinder-cone_01}
local ld = luadraw
local pt3d = ld.pt3d
local M, Mc, O, vecK = pt3d.M, pt3d.Mc, pt3d.Origin,  pt3d.vecK
local sin, cos, pi = math.sin, math.cos, math.pi
local g = ld.graph3d:new{window={-2.5,2.5,-1,7},viewdir={70,70},size={12,12}}
ld.Hiddenlinestyle= "dashed"
local R = 2
local h = 3*R
local H = M(0,0,h)
local u = pt3d.normalize( pt3d.prod(vecK,g.Normal) )
-- points of tangency for the cone
local I, J = pt3d.vecI, pt3d.vecJ -- *(O, I, J)* is an orthogonal coordinate system for the base of the cone.
local f = function(t) 
-- M = O*R*cos(t)*I+R*sin(t)*J is a point on the circular base.
-- f returns the determinant of the vectors g.Normal, M-H, and d(M-H)/dt
    return pt3d.det( g.Normal, O+R*cos(t)*I+R*sin(t)*J-H, -R*sin(t)*I+R*cos(t)*J ) 
end
-- The points of tangency are the points on the circle where the function f vanishes.
local T = ld.solve(f, -pi, pi)
local C, D = Mc(R,T[1],0), Mc(R,T[2],0)
-- 
g:Dcylinder(O,R,H) -- cylinder
g:Dpolyline3d({{C,H,D},{O,H}},"dashed")
g:Ddots3d({O,H})
g:Dlabel3d( "$O$",O,{pos="S"}, "$H$",H,{pos="N"})
g:Show()
\end{luadraw}


\begin{luadraw}{name=cylinder-cone_02}
local ld = luadraw
local pt3d = ld.pt3d
local M, Mc, O,vecK = pt3d.M, pt3d.Mc, pt3d.Origin,  pt3d.vecK
local sin, cos, pi = math.sin, math.cos, math.pi
local g = ld.graph3d:new{window={-2.5,2.5,-1,7},viewdir={70,70},size={12,12}}
ld.Hiddenlinestyle= "dashed"
local R = 2
local h = 3*R
local H = M(0,0,h)
g:Dcone(O,R,H,{color="blue"}) -- cone first
g:Dcylinder(O,R,H,{color="orange", opacity=0.3}) 
g:Ddots3d({O,H})
g:Dlabel3d( "$O$",O,{pos="S"}, "$H$",H,{pos="N"})
g:Dpolyline3d({O,H},"dashed")
g:Show()
\end{luadraw}
\end{document}
